\section{Further progress on the arbitrary-ring case of cote 134-8}

\subsection{F-SM-26: the full interval forces infinite relative Tor amplitude}
Assume the data of BR-13403 and the full interval condition (134.1).
Write \(T=\mathbf1\), and suppose \(k\ne0\).  For every object \(X\), put
\[
 S(X)=\operatorname{coker}\bigl(X\xrightarrow{\alpha_X}J\otimes X\bigr).
\]
Then
\[
 L\varphi(S^nT)\simeq k[n]\qquad(n\ge0).
 \tag{SM26.1}
\]
In the ringoid situation of BR-13405, where
\(\varphi(X)=D\otimes_R X\), this implies
\[
 \operatorname{Tor}^{R}_n(D,S^nT)\cong k\ne0
 \qquad(n\ge0).
 \tag{SM26.2}
\]
Consequently \(D\) has infinite flat dimension, and hence infinite
projective dimension, as a left \(R\)-module.  A finite flat-dimension
hypothesis on \(D\) cannot close the remaining counit test while retaining
the actual nonzero normalized interval hypotheses.

\paragraph{Complete proof.}
The defining sequence
\(0\to X\to J\otimes X\to S(X)\to0\) is exact.
The exact complex-valued representative \(\Lambda\) sends it to a short
exact sequence of complexes.  Its associated triangle, together with
\(\Lambda(J\otimes X)\simeq0\), gives
\[
 L\varphi(SX)\simeq L\varphi(X)[1].
\]
Here homological shift \([1]\) raises the homology degree by one.
Iteration and the canonical normalization \(L\varphi(T)\simeq k[0]\)
give (SM26.1).  Taking degree \(n\) homology gives (SM26.2).
Finite flat dimension \(r\) would make every such Tor group zero for
\(n>r\), a contradiction.  Finite projective dimension implies finite
flat dimension, completing the proof. \(\square\)

\subsection{F-SM-27: an arbitrary-scalar result when \(D\) is of type \(\mathrm{FP}_\infty\)}
Retain the ringoid, scalar, exact-adjunction and full interval hypotheses
of BR-13403--05.  Thus \(R\) is a small \(k\)-linear category, each
\(R(x,y)\) is a projective scalar module,
\(\mathcal M=\operatorname{Fun}_k(R^{\mathrm{op}},\operatorname{Mod}k)\),
\(\varphi=D\otimes_R-\), and the specified adjunction is formed from
a resolution \(B_\bullet\to D\) by coproducts of left representables.
Retain the canonical unit normalization, including degree zero.

Assume the scalar compatibility used in F-SM-24:
\[
 J\otimes(X\otimes_kU)\cong(J\otimes X)\otimes_kU,\qquad
 \alpha_{X\otimes_kU}=\alpha_X\otimes_kU.
 \tag{SM27.1}
\]
This holds for a \(k\)-linear tensor product preserving colimits in each
variable.  Finally suppose that \(D\), as a left \(R\)-module, has a
resolution
\[
 \cdots\longrightarrow A_2\longrightarrow A_1
 \longrightarrow A_0\longrightarrow D\longrightarrow0
 \tag{SM27.2}
\]
in which every \(A_n\) is a finitely generated projective left
\(R\)-module: equivalently, a retract of a finite coproduct of left
representables.  This is the additional condition
\(D\in\mathrm{FP}_\infty(R)\); no finite length for (SM27.2) is required.

Then the exact adjunction specified in BR-13405 induces
\[
 W^{-1}\mathcal M\simeq D_{\ge0}(k).
 \tag{SM27.3}
\]
There is no cardinality, ordinary injective-dimension or pure
injective-dimension restriction on \(k\), and the scalar modules
\(R(x,y)\) need not be finitely generated.  More strongly, with
\[
 C_x(U)(y)=\operatorname{Hom}_k(R(x,y),U),
 \tag{SM27.4}
\]
\(L\varphi\) vanishes on every \(C_x(U)\) and on every product of such
objects, for arbitrary scalar modules \(U\).

\paragraph{Exact known input.}
We use Benson and Goodearl's (2000) periodic-flat-module theorem:
a flat module \(V\) in an exact sequence
\(0\to V\to P\to V\to0\), with \(P\) projective, is projective.
The precise primary locator is Theorem 2.5, printed p.~50
(also stated as Theorem 1.1 on p.~46), in
\emph{Periodic flat modules, and flat modules for finite groups}.
Its proof occupies pp.~50--51, after Lemmas 2.1--2.4.
The complete deduction from that known theorem to the contractibility
needed below is given here.  In particular, ordinary acyclicity of a
complex of projectives, unbounded to the right cohomologically, is
not substituted for pure acyclicity.

\paragraph{Complete proof.}
We use the convention that a pure monomorphism of scalar modules
remains injective after tensoring with every scalar module.
A scalar module \(E\) is pure-injective if maps to \(E\) extend across
all such monomorphisms.  For every scalar module \(V\), its character
module
\[
 V^+=\operatorname{Hom}_{\mathbf Z}(V,\mathbf Q/\mathbf Z)
\]
is pure-injective: tensor--Hom adjunction identifies extension across
a pure monomorphism with extension to the injective abelian group
\(\mathbf Q/\mathbf Z\).  We will use this elementary fact both to
produce tests and to detect purity.

\emph{1. Coinduction with pure-injective coefficients is annihilated.}
There is a natural adjunction
\[
 \operatorname{Hom}_{\mathcal M}(X,C_x(E))
   =\operatorname{Hom}_k(X(x),E).
 \tag{SM27.5}
\]
Consequently, when \(E\) is pure-injective, \(C_x(E)\) is injective
with respect to all monomorphisms that are scalar-pure at every
object of \(R\).  Condition (SM27.1) and the interval's monomorphism
condition show that \(\alpha_X:X\to J\otimes X\) is objectwise
scalar-pure: tensoring it with any \(U\) gives the monomorphism
\(\alpha_{X\otimes_k U}\).  Thus \(\alpha_{C_x(E)}\) splits in
\(\mathcal M\), by extending the identity using (SM27.5).
The interval condition gives
\[
 L\varphi(C_x(E))=0\qquad(E\text{ scalar pure-injective}).
 \tag{SM27.6}
\]
This argument does not declare \(C_x(E)\) ordinarily injective.

\emph{2. Dualizing the finite-projective resolution.}
Fix \(x\), and form the cochain complex
\[
 K_x^n=\operatorname{Hom}_{R\text{-}\mathrm{Mod}}
                (A_n,R(x,-)),\qquad
 \delta^n(f)=f\circ d_{n+1}.
 \tag{SM27.7}
\]
Every \(K_x^n\) is a projective scalar module: for
\(A_n=R(y,-)\) it is \(R(x,y)\), and the general case follows by
finite coproducts and retracts.  For every scalar module \(U\),
evaluation and finite-projective duality give natural isomorphisms
\[
 A_n\otimes_R C_x(U)
   \cong\operatorname{Hom}_k(K_x^n,U).
 \tag{SM27.8}
\]
For a representable this is the defining formula (SM27.4).
Finite coproducts convert to finite products, which coincide, and
retracts preserve the isomorphism.  This also proves naturality in
\(A_n\); hence the isomorphisms commute with the differentials.
The homological complex \(A_\bullet\otimes_R C_x(U)\), computing
\(L\varphi(C_x(U))\), is therefore the transpose complex
\(\operatorname{Hom}_k(K_x^\bullet,U)\).

By (SM27.6), this transpose complex is acyclic for every pure-injective
\(U\).  In particular take the injective cogenerator
\(E_0=\operatorname{Hom}_{\mathbf Z}(k,\mathbf Q/\mathbf Z)\).
Then
\(\operatorname{Hom}_k(K_x^\bullet,E_0)
 =\operatorname{Hom}_{\mathbf Z}(K_x^\bullet,\mathbf Q/\mathbf Z)\);
exactness and the cogenerator property imply that \(K_x^\bullet\)
is acyclic.

Let \(Z^n=\ker(\delta^n)\).  As \(K_x^\bullet\) starts in degree zero,
acyclicity implies \(Z^0=0\) and gives exact sequences
\[
 0\longrightarrow Z^n\longrightarrow K_x^n
   \xrightarrow{\delta^n}Z^{n+1}\longrightarrow0.
 \tag{SM27.9}
\]
These sequences are pure.  To see this directly, fix a pure-injective
\(E\) and a map \(f:Z^n\to E\), with \(n\ge1\).
The map \(f\delta^{n-1}:K_x^{n-1}\to E\) is a cycle in the
transpose complex.  Its acyclicity gives \(g:K_x^n\to E\) with
\(g\delta^{n-1}=f\delta^{n-1}\).  Surjectivity onto \(Z^n\)
implies \(g|_{Z^n}=f\).  The assertion for \(n=0\) is immediate.
In particular this extension property holds for every \(E=V^+\).
Tensor--Hom adjunction and the fact that
\(\mathbf Q/\mathbf Z\) detects nonzero elements imply that
\(V\otimes_k Z^n\to V\otimes_k K_x^n\) is injective for every
\(V\).  This proves purity of (SM27.9).

\emph{3. Contractibility, using the genuine periodicity input.}
A pure submodule of a flat module is flat.  In the present instance
one can check this without another theorem: for an injection
\(V'\to V\), put \(Z^n\to K_x^n\) in the tensor square.
Both horizontal maps are injective by purity, and the right vertical
map is injective by flatness of \(K_x^n\).  Therefore the left
vertical map is injective.  This is the flatness criterion for
\(Z^n\).

Take the direct sum of (SM27.9) for \(n\ge0\).
Since \(Z^0=0\), the two end sums are canonically isomorphic after
shifting their indices.  We obtain
\[
 0\longrightarrow Z\longrightarrow P\longrightarrow Z
 \longrightarrow0,\qquad
 Z=\bigoplus_{n\ge0}Z^n,\quad P=\bigoplus_{n\ge0}K_x^n.
 \tag{SM27.10}
\]
Here \(Z\) is flat and \(P\) is projective.
The cited Benson--Goodearl Theorem 2.5 makes \(Z\) projective.
Each \(Z^n\), as a direct summand of \(Z\), is consequently
projective.  Thus every sequence (SM27.9) splits.
Choosing splittings identifies
\(K_x^n=Z^n\oplus Z^{n+1}\) and
\(\delta^n(z,w)=(w,0)\); the maps
\(h^n(z,w)=(0,z)\) give a contracting homotopy
(with \(h^0=0\)).  This proves contractibility of \(K_x^\bullet\),
including its unbounded right-hand end.

For every scalar module \(U\), transpose this homotopy.
Then \(\operatorname{Hom}_k(K_x^\bullet,U)\) is contractible.
By (SM27.8),
\[
 L\varphi(C_x(U))=0\qquad\text{for every }x,U.
 \tag{SM27.11}
\]

\emph{4. Products and the specified adjunction.}
Each functor \(A_n\otimes_R-\) commutes with arbitrary products,
because \(A_n\) is a retract of a finite sum of representables,
whose tensor functors are evaluation functors.
Therefore for any family \(X_i\) there is a natural isomorphism of
complexes
\[
 A_\bullet\otimes_R\prod_iX_i
   \cong\prod_i(A_\bullet\otimes_R X_i).
 \tag{SM27.12}
\]
Products of scalar modules are exact, so a product of acyclic
complexes is acyclic.  Formula (SM27.11) now proves the stronger
vanishing assertion for every product of coinduced objects.
Notice that this product argument computes \(L\varphi\) with \(A\);
it does not assume that the originally chosen \(B_n\) is finite.

Let \(P\) be an acyclic projective object of
\(\operatorname{Ch}_{\ge0}(k)\).  BR-13402 writes
\[
 P=\bigoplus_{n\ge0}D^{n+1}(Q_n)
   =\prod_{n\ge0}D^{n+1}(Q_n),
 \tag{SM27.13}
\]
where each \(Q_n\) is projective; the equality of sum and product is
degreewise, since only two summands occur in any chain degree.
The right adjoint \(G\) preserves products.  For every \(y\),
\[
 G(D^{n+1}Q_n)(y)
   =\operatorname{Hom}_{\mathrm{Ch}}(B_\bullet(y),D^{n+1}Q_n)
   =\operatorname{Hom}_k(B_n(y),Q_n).
\]
The latter formula is the disk adjunction, obtained by sending the
degree-\(n\) map freely to \(Q_n\); its degree-\(n+1\) companion is
determined by the differential.  Write the chosen resolution term as
\(B_n=\bigoplus_{a\in I_n}R(x_{n,a},-)\).
It follows naturally as right \(R\)-modules that
\[
 GP\cong\prod_{n\ge0}\prod_{a\in I_n}C_{x_{n,a}}(Q_n).
 \tag{SM27.14}
\]
The preceding product vanishing gives \(L\varphi(GP)=0\).
BR-13403 and F-SM-17 conclude (SM27.3).  This last invocation uses
the canonical unit normalization to handle projective complexes'
degree-zero summands; the coinduction vanishing itself did not need
that normalization. \(\square\)

\paragraph{Exact remaining scope.}
Condition \(\mathrm{FP}_\infty\) is an added hypothesis on \(D\),
not a consequence asserted here of the existence of finitely
presented projective generators of \(\mathcal M\).
It supplies finite-projective duality in (SM27.8) and product
commutation in (SM27.12).  Neither identity has been asserted for an
arbitrary infinite coproduct of representables.
Thus F-SM-27 adds a third proved range to F-SM-24:
arbitrary scalar rings and arbitrary projective scalar morphism
modules, provided \(D\) has this finite-generation property in every
resolution degree.  The case without any of these extra hypotheses
remains the exact test (134.6).

A countable-subring argument has not removed that last condition.
Although a tensor element uses finitely many tensors, its coefficients
in \(GP\) are global linear maps on the modules \(B_n(y)\).
A boundary found after restriction to a smaller ringoid can require
extending such maps back to all of \(B_n(y)\), compatibly with the
ringoid action.  Finite support of the outer tensor expression does
not provide that extension.  No automatic descent or compatible
extension assertion is used in the proved result.

\paragraph{Primary reference.}
Benson, D. J., \& Goodearl, K. R. (2000).
Periodic flat modules, and flat modules for finite groups.
\emph{Pacific Journal of Mathematics, 196}(1), 45--67.
\url{https://doi.org/10.2140/pjm.2000.196.45}.
Theorem 2.5, p.~50; proof pp.~50--51.
Primary publisher PDF:
\url{https://msp.org/pjm/2000/196-1/pjm-v196-n1-p03-s.pdf}.

